\chapter{Collaboration and Reproducibility}
\label{chap:collaboration}
\index{collaboration}\index{reproducibility}\index{version control}\index{Git}\index{backup}\index{README file}

A document that builds only in one author's long-lived folder is not yet a
reproducible scholarly object. Collaboration exposes hidden assumptions: an
unshared image, a local font, an absolute path, an unexplained command, or a
revision whose author and purpose cannot be recovered.

\begin{learningoutcomes}
  \item Organise and name project files for readers and collaborators.
  \item Write comments, semantic commands, and a practical README.
  \item Compare versioned copies, cloud editing, change marking, and Git.
  \item Share a complete, licensed, permission-aware source package.
  \item Reproduce a document build in a clean environment.
  \item Design backups and releases that can be restored and audited.
\end{learningoutcomes}

\section{Reproducible means reconstructable}

For a \LaTeX{} document, reproducible compilation means that an authorised reader
can obtain the stated output from the preserved source, dependencies, assets,
and instructions. It does not imply that the underlying scientific result is
valid. Computational reproducibility additionally asks whether data and code
regenerate the reported values, tables, and figures.

Record the boundary. A document may be fully rebuildable while sensitive raw
data remain restricted. In that case, provide permitted derived material,
metadata, access conditions, and a test data set where useful.

\begin{classroomnote}
``It compiles on my computer'' is evidence about one environment. A clean build
from a newly copied project tests whether the project itself contains the
dependencies it claims to contain.
\end{classroomnote}

\section{Names are navigation}

Choose short, descriptive, portable names:

\begin{terminalcode}{A readable project}
measurement_article/
|-- main.tex
|-- preamble.tex
|-- sections/methods.tex
|-- figures/method_comparison.pdf
|-- data/simulated_measurements.csv
|-- analysis/check_results.R
|-- references.bib
`-- README.md
\end{terminalcode}

Prefer one case convention and avoid ambiguous dates. ISO-style dates such as
\file{2026-08-01} sort chronologically. Spaces can be supported, but simple
ASCII names with hyphens or underscores reduce shell, archive, and service
friction. Never encode confidential participant information in filenames.

Generated files need not masquerade as source. Keep build products in a
documented output directory or ignore them in version control. Do not call a
working draft \file{final-final-2.tex}; use a revision system and meaningful
release tags.

\begin{tryit}
List the ten most important files in one current project. Rename two ambiguous
files on a disposable copy, update their references, and compile. Note whether
case-only changes behave correctly on your file system.
\end{tryit}

\section{Comments explain reasons}

The percent sign begins a \LaTeX{} comment. Use comments for purpose, constraints,
and temporary diagnostics:

\begin{latexcode}{Comments that retain context}
% Journal requires continuous line numbers for review.
% Remove the review option only after editorial acceptance.
\documentclass[review]{publisherclass}

% Equality line, not a fitted regression.
\includegraphics{figures/method_comparison.pdf}
\end{latexcode}

Avoid comments that repeat the code or preserve whole obsolete paragraphs
indefinitely. Version history can recover removed text. Delete resolved draft
notes before release, and search for project-specific markers such as
\code{TODO}, \code{FIXME}, and \code{[PLACEHOLDER]}.

\section{Custom commands define a shared language}

Semantic commands give collaborators one controlled notation:

\begin{latexcode}{Shared scientific notation}
\newcommand{\MeanDifference}{\ensuremath{\bar d}}
\newcommand{\MethodA}{\textsc{Method A}}
\end{latexcode}

Place them in a small shared preamble file and document non-obvious arguments.
Choose names that reveal meaning. Do not redefine familiar \LaTeX{} commands or
create private abbreviations for ordinary prose.

\begin{pausepredict}
If the notation for the mean difference changes from $\bar d$ to $\Delta$, how
many edits are required when the source uses a semantic command? What new risk
appears if some authors typed the symbol directly?
\end{pausepredict}

\section{Four ways to manage revision}

Different teams need different levels of history and access.

\subsection{Numbered snapshots}

For a brief solo project, dated archive copies can preserve milestones. Keep one
working source and write protected snapshots outside it. A directory containing
five competing main files invites compilation of the wrong version.

\subsection{Cloud \LaTeX{} collaboration}

Services such as Overleaf provide shared projects, an editor, compilation,
comments, and history features whose exact availability can depend on the
current service and plan. Agree on the main file, TeX engine, review practice,
and who resolves conflicts. Export complete source snapshots at milestones;
cloud convenience is not a backup policy by itself.

\subsection{Visible change marking}

Packages such as \pkg{changes} or \pkg{latexdiff} can show additions and
deletions. Their output is a review artefact, not automatically the final source.
Complex commands, moved text, citations, and floats need manual inspection.
Retain an accepted clean version and verify that no change markup survives.

\subsection{Git and hosted repositories}

Git records snapshots of text files efficiently. Hosting services add remote
collaboration, issues, review, and releases. This is optional advanced tooling,
not a prerequisite for writing \LaTeX{}.

\begin{terminalcode}{A small Git cycle}
git status
git diff
git add chapters/methods.tex
git commit -m "Clarify paired-difference definition"
\end{terminalcode}

Commits should describe a coherent change. Inspect the diff before staging.
Never commit credentials, private keys, confidential raw data, or files whose
licence prohibits redistribution. A public repository is a publication action.

\begin{advancedtopic}
A useful branch strategy can be modest: one protected main branch, short-lived
branches for substantial changes, review before merging, and tagged releases
for submissions. Complex branching without a team agreement can obscure rather
than improve the scholarly record.
\end{advancedtopic}

\section{Resolve collaboration at the level of ownership}

Assign responsibility by scholarly unit: one author owns Methods, another
validates analyses, and a corresponding author controls the submission release.
Ownership does not prevent review; it prevents simultaneous uncoordinated edits
to the same sentences.

Record decisions in an issue, meeting note, or commit message. A useful record
states the question, evidence considered, decision, owner, and date. Email alone
can scatter the rationale across private inboxes.

\begin{commonerror}
\textbf{Symptom:} two correct revisions overwrite one another or leave half of
each argument in the final paragraph.

\textbf{Cause:} both collaborators edited the same unit without a merge and
review convention.

\textbf{Repair:} designate an owner, compare changes explicitly, resolve the
scientific meaning before resolving punctuation, and compile the merged result.
\end{commonerror}

\section{README as the first reproducibility test}

A README should answer:

\begin{itemize}
  \item What is this project, and are its data real or simulated?
  \item Which file is the official entry point?
  \item Which TeX engine, bibliography tool, and extra tools are required?
  \item What command creates the document and generated figures or tables?
  \item Where do inputs come from and what are their licences or restrictions?
  \item Which placeholders or manual steps remain?
  \item Where should generated output appear?
\end{itemize}

\begin{latexcode}{A concise build statement}
% README.md (shown as plain text)
Build: latexmk -pdf main.tex
Main file: main.tex
Data: simulated teaching values in data/
Analysis check: Rscript analysis/check_results.R
\end{latexcode}

The listing is README text, not commands to place in a TeX source file. Use
Markdown code fences in the actual README.

\section{Pin what affects the build}

Record the engine and important tool versions. A minimal environment record can
include TeX distribution and date, \pkg{latexmk}, Biber or BibTeX, custom class
versions, and required fonts. If a package behaviour is critical, preserve the
authorised class or style file and its licence when distribution permits.

A container or continuous-integration build can strengthen repeatability, but
it also needs maintenance. The simplest adequate solution may be a current TeX
Live release plus a documented command. Do not introduce a complex platform
that no collaborator can repair.

\section{Avoid absolute paths}

This source works only for one account:

\begin{latexcode}{Non-portable figure path}
\includegraphics{/Users/name/Desktop/project/figure1.pdf}
\end{latexcode}

Use a project-relative path:

\begin{latexcode}{Portable figure path}
\includegraphics{figures/figure1.pdf}
\end{latexcode}

The same rule applies to bibliography databases, data imports, and scripts.
Check case exactly: \file{Figure1.pdf} and \file{figure1.pdf} may be different
files on the collaborator's system.

\section{Backups need independent copies}

Synchronisation protects availability but can synchronise accidental deletion
or corruption. A robust pattern keeps multiple copies, on more than one medium,
with at least one versioned or offline/off-site copy. Test restoration.

Back up source, data permitted by governance, analysis, figure sources,
bibliography, README, licences, approvals, and released PDFs. Auxiliary files
can usually be rebuilt. Encrypt and control access where required; a backup is
not exempt from data-protection obligations.

\section{Licence and permission awareness}

Distinguish four questions:

\begin{enumerate}
  \item May the team use the material for this research?
  \item May it be redistributed with source?
  \item May it be modified?
  \item What attribution, notice, or share-alike terms apply?
\end{enumerate}

Copyright, data licences, software licences, consent, confidentiality, and
publisher agreements are related but not interchangeable. Record the source,
creator, licence, permission, and required credit for each external asset.

\begin{scienceinpractice}
A geological report may use an openly licensed base map, confidential sampling
coordinates, an author's original plot, and a publisher-reused diagram under
written permission. The source package cannot simply be made public as one
undifferentiated object.
\end{scienceinpractice}

\section{The clean-room rehearsal}

Create a fresh directory, copy only the files named by the release manifest,
and follow the README as if you were a new collaborator. A successful rehearsal
checks filenames, paths, tool order, missing dependencies, and unstated manual
steps. Compare the resulting page count or checksum only when deterministic
output is expected; timestamps and tool versions can legitimately change bytes.

\begin{goodpractice}
Make the release from a known version, run the numerical verification, build
the document, inspect the log and PDF, and archive the exact inputs and outputs
together. A release is more than the latest file in a synchronised folder.
\end{goodpractice}

\chapterproject

\begin{miniproject}
Open the running article's README. From a fresh copy, run the base-R result check
and compile \file{main.tex}. Add one comment explaining why the difference is
defined as B minus A and one semantic command for a repeated notation. Create a
manifest of the files another scientist needs; exclude rebuildable auxiliary
files and identify the original source of each figure.
\end{miniproject}

\chaptersummary

Collaboration depends on shared meaning, visible ownership, and reconstructable
releases. Descriptive names, purposeful comments, semantic commands, one main
file, a complete README, relative paths, and tested backups reduce hidden
assumptions. Cloud editors, change-marking tools, and Git solve different
revision problems. Reproducible sharing also requires licences, permissions,
access controls, environment records, a clean build, and an exact release
manifest.

\chaptercheck
\begin{chapterchecklist}
  \item File and folder names are descriptive, portable, and privacy-safe.
  \item Comments explain constraints and no unresolved draft markers remain.
  \item The main file, engine, tools, and build commands are documented.
  \item Collaborators know ownership, review, and merge conventions.
  \item Paths are relative and all permitted dependencies are included.
  \item Licences, permissions, and restricted materials are recorded.
  \item Backups can be restored and the release builds in a clean directory.
\end{chapterchecklist}

\chapterexercisestitle
\begin{chapterexercises}
  \item Rename four poorly named files and explain your convention.
  \item Distinguish a synchronised folder, version history, and a backup.
  \item Write a README for a three-file article and its bibliography.
  \item Replace two absolute paths with project-relative paths.
  \item Compare numbered snapshots, change markup, Overleaf, and Git for a team.
  \item Design an asset-permission register for four figure types.
  \item Plan a clean-room build test and define its success evidence.
  \item Challenge: create a release manifest that separates public, restricted,
        generated, and rebuildable files.
\end{chapterexercises}
